\section{Related Work}

\paragraph{Chart question answering.}
Chart question answering requires models to combine visual perception, OCR, structural grounding, and numerical reasoning. Earlier benchmarks such as DVQA, FigureQA-style synthetic chart reasoning, and PlotQA emphasize controlled visual reasoning over chart data \citep{kafle2018dvqa,methani2020plotqa}. ChartQA \citep{masry2022chartqa} evaluates both human-written and machine-generated chart questions that involve value lookup, comparison, and arithmetic. Recent chart-specialized systems add plot-to-table reasoning, chart derendering, chart-specific pretraining, and visual instruction tuning \citep{liu2023deplot,liu2023matcha,masry2023unichart,han2023chartllama,masry2024chartgemma}. Recent benchmarks stress more complex and realistic charts \citep{chartbench2024,wang2024charxiv,masry2025chartqapro,wei2025chartmind}, and ChartInsights \citep{wu2024chartinsights} evaluates low-level chart tasks with fine-grained error taxonomies. These diagnostics slice errors by task type, chart type, or reasoning step; none isolates the case in which the model finds the right quantity but reports it under the wrong scale, nor tracks the \emph{direction} of that error. ScaleQA-Natural is complementary: it is a narrow, directional slice that can be mined from any of these benchmarks, and our fix--break accounting applies to any correction evaluated on them.

\paragraph{Numerical reliability in MLLMs.}
Numerical errors in multimodal models can arise from visual grounding failures, OCR mistakes, wrong axis interpretation, arithmetic errors, or answer normalization artifacts. Scale inconsistency differs from a generic numeric miss because the prediction is close under a simple multiplicative transformation. This makes it easy to miss in aggregate error analysis and risky in deployment, since an answer like 0.61 can look superficially related to 61 while being semantically wrong. The distinction also changes the evaluation target: the question is not only whether a model can compute or retrieve a number, but whether it can choose the output convention supported by the chart's axis, unit, and question wording.

\paragraph{Diagnostic evaluation and correction policies.}
Most chart QA work reports aggregate task accuracy, while reliability failures are often analyzed post hoc. We instead treat scale conflicts as a diagnostic slice with inclusion criteria, an audit trail, and fix--break accounting. A useful correction must recover scale errors, preserve non-scale examples, and avoid applying the wrong direction. The high oracle gaps in our experiments further distinguish candidate generation from candidate selection, a failure mode not captured by single-answer accuracy alone. We use TextVQA \citep{singh2019towards} only as an adjacent text-rich safety check.

\paragraph{Parameter-efficient and preference-style adaptation.}
LoRA updates a model through low-rank trainable matrices while keeping the base weights fixed \citep{hu2021lora}. This makes it attractive for targeted reliability interventions in large MLLMs. Contrastive and preference-style objectives, including direct preference optimization \citep{rafailov2023direct}, optimize relative preferences between candidate outputs. Our objective is more local: for a chart and question, we contrast a preferred numeric answer with a scale-related alternative, and study how the direction of this contrast changes behavior.
